%% kapa: exam=1 round=37 year=2026 period=1 form=A subject=economics subject_ko=경제학원론 seq=6 number=46
%% source: Q-Net 감정평가사 기출문제 (cst00309) artlSeq=5258924 file=raw/1008/5258924/2026년 제37회 감정평가사 1차 1교시.pdf
%% conversion: kordoc 4.24.51 -> tools/gen1.py (auto); review=proofread-sonnet; flags=none
\begin{question}{46}
하이브리드 자동차 시장을 양분하고 있는 H사와 T사의 광고여부에 따른 보수행렬은 아래와 같다. 다음 중 옳은 것을 모두 고른 것은? (단, 각 보수쌍에서 왼쪽은 H사의 보수이고, 오른쪽은 T사의 보수이다.)
\par
{\small\centering
\begin{tabular}{|c|c|c|c|}\hline
\multirow{4}{*}{H사} & \multirow{2}{*}{~} & \multicolumn{2}{c|}{T사} \\ \cline{3-4}
 & & 광고함 & 광고안함 \\ \cline{2-4}
 & 광고함 & (2500, 1500) & (3000, 0) \\ \cline{2-4}
 & 광고안함 & (1500, 2000) & (4000, 500) \\ \hline
\end{tabular}\par}
\begin{framedtext}
ㄱ. H사는 우월전략을 가지고 있다. \par
ㄴ. T사의 우월전략은 광고를 하는 것이다. \par
ㄷ. 내쉬균형은 두 회사 모두 광고를 하는 것이다. \par
ㄹ. 이 게임의 우월전략균형은 두 회사 모두 광고를 안하는 것이다. \par
\end{framedtext}
\begin{choices}
  \item ㄱ, ㄴ
  \item ㄴ, ㄷ
  \item ㄷ, ㄹ
  \item ㄱ, ㄴ, ㄹ
  \item ㄴ, ㄷ, ㄹ
\end{choices}
\end{question}
\answer{2}
